\section{Instrumentos y Equipos}
\begin{itemize}
    \item Amperímetro
    \item Voltímetro
    \item Vatímetro
    \item Osciloscopio
    \item Fuente de tensión alterna ajustable (Variac)
    \item Multímetro digital
    \item Puente rectificador de onda completa
    \item Resistencia de \SI{33}{\ohm} a \SI{25}{\celsius}
    \item Condensador de \SI{33}{\micro\farad}
    \item Cables de conexión
\end{itemize}

\section{Procedimiento de Laboratorio}

\begin{figure}[H]
\centering
\begin{minipage}{0.48\textwidth}
\centering
\begin{circuitikz}[american, scale=0.6, transform shape]
    \draw
    %% Fuente de CA variable (Variac)
    (0,0) to[sV, l_={120\,V, 60\,Hz}] (0,2)
    
    %% Flecha de ajuste variable sobre el generador
    (0.35,0.2) -- (-0.35,1.8)
    
    %% AC+ → nodo izquierdo del puente (Iac sobre el positivo)
    (0,2) to[short, i=$I_{AC}$] (2.5,2)
    
    %% AC- → nodo derecho del puente (pasa por debajo del diamante)
    (0,0) -- (5.5,0) -- (5.5,2)
    
    %% Puente rectificador (rombo de lados iguales)
    %% Centro (4,2), Izq (2.5,2), Der (5.5,2), Sup (4,3.5), Inf (4,0.5)
    (2.5,2) to[D, l=D$_1$] (4,3.5)
    (4,0.5) to[D, l=D$_4$] (2.5,2)
    (5.5,2) to[D, l_=D$_2$] (4,3.5)
    (4,0.5) to[D, l_=D$_3$] (5.5,2)
    
    %% Salida DC
    (4,3.5) -- (7,3.5)
    (4,0.5) -- (7,0.5)
    
    %% Carga: solo resistencia
    (7,3.5) to[R, l=$33\,\Omega$, i=$I_{DC}$] (7,0.5)
    
    %% Extensión para voltímetro a la derecha del circuito
    (7,3.5) -- (9.5,3.5)
    (7,0.5) -- (9.5,0.5)
    %% Voltímetro
    (9.5,3.5) to[open, v=$V_{DC}$] (9.5,0.5)
    ;
\end{circuitikz}
\vspace{0.3cm}
\par (a) Montaje para ondas sinusoidales rectificadas
\end{minipage}
\hfill
\begin{minipage}{0.48\textwidth}
\centering
\begin{circuitikz}[american, scale=0.6, transform shape]
    \draw
    %% Fuente de CA variable (Variac)
    (0,0) to[sV, l_={120\,V, 60\,Hz}] (0,2)
    
    %% Flecha de ajuste variable sobre el generador
    (0.35,0.2) -- (-0.35,1.8)
    
    %% AC+ → nodo izquierdo del puente (Iac sobre el positivo)
    (0,2) to[short, i=$I_{AC}$] (2.5,2)
    
    %% AC- → nodo derecho del puente (pasa por debajo del diamante)
    (0,0) -- (5.5,0) -- (5.5,2)
    
    %% Puente rectificador (rombo de lados iguales)
    (2.5,2) to[D, l=D$_1$] (4,3.5)
    (4,0.5) to[D, l=D$_4$] (2.5,2)
    (5.5,2) to[D, l_=D$_2$] (4,3.5)
    (4,0.5) to[D, l_=D$_3$] (5.5,2)
    
    %% Salida DC
    (4,3.5) -- (6.5,3.5)
    (4,0.5) -- (7,0.5)
    
    %% Amperímetro en línea común DC+
    (6.5,3.5) to[short, i=$I_{DC}$] (7,3.5)
    
    %% Carga: Resistencia y Capacitor en paralelo
    (7,3.5) to[R, l=$33\,\Omega$] (7,0.5)
    (7,3.5) -- (8.5,3.5)
    (8.5,3.5) to[C, l=$33\,\mu$F] (8.5,0.5)
    (8.5,0.5) -- (7,0.5)
    
    %% Extensión para voltímetro a la derecha del circuito
    (8.5,3.5) -- (10.5,3.5)
    (8.5,0.5) -- (10.5,0.5)
    %% Voltímetro
    (10.5,3.5) to[open, v=$V_{DC}$] (10.5,0.5)
    ;
\end{circuitikz}
\vspace{0.3cm}
\par (b) Montaje para ondas no sinusoidales
\end{minipage}
\caption{Diagramas de montaje del procedimiento experimental.}
\label{fig:montajes-experimentales}
\end{figure}

\subsection{Medición de ondas sinusoidales o sinusoidales rectificadas}
\begin{enumerate}
    \item Realizar un montaje con una fuente alterna ajustable (fuente de \SI{120}{V} a \SI{60}{Hz} en serie con un autotransformador o Variac) que alimente una carga compuesta por un puente rectificador de onda completa, el cual alimenta en corriente continua un elemento resistivo de \SI{33}{\ohm} a \SI{25}{\celsius}.
    \item Medir con instrumentos de los tipos descritos en el marco teórico (bobina móvil, hierro móvil y electrodinámico) los valores de tensión y corriente del lado de corriente continua, así como la corriente del lado de corriente alterna.
    \item Calcular el coeficiente de conversión que introduce un equipo magnetoeléctrico para dar valores eficaces. Comprobar este valor calculado con las mediciones realizadas.
\end{enumerate}

\subsection{Medición de ondas no sinusoidales}
\begin{enumerate}
    \item Realizar un montaje con una fuente alterna ajustable (fuente de \SI{120}{V} a \SI{60}{Hz} en serie con un autotransformador o Variac) que alimente una carga compuesta por un puente rectificador de onda completa, el cual alimenta en corriente continua un elemento resistivo de \SI{33}{\ohm} en paralelo con un condensador de aproximadamente \SI{33}{\micro\farad}.
    \item Medir con instrumentos de los tipos descritos en el marco teórico los valores de tensión y corriente del lado de corriente continua y la corriente del lado de corriente alterna.
    \item Calcular el coeficiente de conversión que introduce un equipo magnetoeléctrico para dar valores eficaces. Comprobar este valor calculado con las mediciones realizadas.
    \item Realizar varias mediciones para tensiones de alimentación diferentes y deducir el valor de la resistencia utilizada.
    \item Trazar la curva de tensión versus corriente sobre papel milimetrado, marcando cada punto de medición con sus correspondientes errores.
    \item Explicar las discrepancias entre las mediciones del punto~1 y las realizadas en el punto~2, utilizando un osciloscopio para observar las formas de onda de la corriente y la tensión.
\end{enumerate}
